This chapter establishes the experimental framework and software architecture used to empirically evaluate the performance, operational cost, and output quality of Retrieval-Augmented Generation (RAG) optimization techniques. To address the research questions formulated in Chapter~\ref{sec:introduction}, this study implements a quantitative benchmarking suite built upon an enterprise-ready Java and Spring AI infrastructure connected to Google Cloud Platform (GCP) services. The remainder of this chapter is structured as follows: Section~\ref{subsec:res_methodology} details the overall research approach and experimental design, while Section~\ref{subsec:sys_implementation} describes the software architecture and cloud infrastructure. Section~\ref{subsec:opt_pipeline} outlines the specific optimization strategies evaluated. Finally, Section~\ref{subsec:bench_datasets} introduces the target datasets and Section~\ref{subsec:eval_metrics} specifies the quantitative evaluation metrics and measurement instrumentation.

\subsection{Research Methodology}
\label{subsec:res_methodology}

Empirical research is conducted using quantitative benchmarking in a two-stage experimental strategy. In the first phase, each of the optimization methods is tested against the baseline solution. In the second phase, pairwise and multi-method combinations are tested until the full-stack solution is tested against the baseline to identify the Pareto frontier of the methods and study the potential systemic trade-offs between optimization methods. The combination matrix is presented in Table~\ref{tab:combination_matrix}. The target of this research is to address the potential benefits of addressing the research questions outlined in Table~\ref{tab:research_questions}. 10 iterations per configuration are executed to test statistical significance and to mitigate potential network jitters and cold-start biases, and to make sure that caching is utilized effectively. 

\begin{table}[htbp]
\centering
\caption{Full factorial experimental design matrix ($2^4 = 16$ configurations) across the two-stage benchmarking framework.}
\label{tab:combination_matrix}
\begin{tabular}{l l c c c c}
\toprule
\textbf{Config} & \textbf{Pipeline Composition} & \textbf{$M_1$: S2B} &\textbf{$M_2$: SR} & \textbf{$M_3$: PC} &
\textbf{$M_4$: MC} \\
    \midrule
\multicolumn{6}{l}{\textbf{Baseline ($k=0$)}} \\
    $C_{0}$  & Naive RAG Baseline            & --           & --          & --          & --          \\
    \midrule
\multicolumn{6}{l}{\textbf{Stage 1: Single-Factor Isolated Testing ($k=1$)}} \\
    $C_{1}$  & Small-to-Big Retrieval (S2B)   & \checkmark   & --          & --          & --          \\
    $C_{2}$  & Semantic Re-ranking (SR)       & --           & \checkmark  & --          & --          \\
    $C_{3}$  & Prompt Caching (PC)            & --           & --          & \checkmark  & --          \\
    $C_{4}$  & Model Cascading (MC)           & --           & --          & --          & \checkmark  \\
    \midrule
\multicolumn{6}{l}{\textbf{Stage 2: Pairwise Combinations ($k=2$)}} \\
    $C_{5}$  & S2B + SR                       & \checkmark   & \checkmark  & --          & --          \\
    $C_{6}$  & S2B + PC                       & \checkmark   & --          & \checkmark  & --          \\
    $C_{7}$  & S2B + MC                       & \checkmark   & --          & --          & \checkmark  \\
    $C_{8}$  & SR + PC                        & --           & \checkmark  & \checkmark  & --          \\
    $C_{9}$  & SR + MC                        & --           & \checkmark  & --          & \checkmark  \\
    $C_{10}$ & PC + MC                        & --           & --          & \checkmark  & \checkmark  \\
    \midrule
\multicolumn{6}{l}{\textbf{Stage 2: Multi-Module Triplets ($k=3$)}} \\
    $C_{11}$ & S2B + SR + PC                  & \checkmark   & \checkmark  & \checkmark  & --          \\
    $C_{12}$ & S2B + SR + MC                  & \checkmark   & \checkmark  & --          & \checkmark  \\
    $C_{13}$ & S2B + PC + MC                  & \checkmark   & --          &\checkmark  & \checkmark  \\
    $C_{14}$ & SR + PC + MC                   & --           & \checkmark  &\checkmark  & \checkmark  \\
\midrule
\multicolumn{6}{l}{\textbf{Stage 2: Full-Stack Integration ($k=4$)}} \\
    $C_{15}$ & S2B + SR + PC + MC (Full-Stack)& \checkmark   & \checkmark  &\checkmark  & \checkmark  \\
\bottomrule
\end{tabular}

\vspace{0.75em}
\parbox{0.88\textwidth}{\footnotesize
\textbf{Notes:}
$M_1$: Small-to-Big Retrieval (S2B); $M_2$: Semantic Re-ranking (SR); $M_3$: Prompt Caching (PC); $M_4$: Model Cascading (MC). Checkmarks (\checkmark) denote active optimization components in the pipeline.}
\end{table}
% DRAFT NOTES:
% - Experimental quantitative benchmarking approach.
% - Addressing RQ1 (Latency/TTFT), RQ2 (Cost $), RQ3 (Quality/RAGAS), RQ4 (Pareto Equilibrium).
% - Multi-iteration execution (5 iterations per configuration) to mitigate network jitter and cold-start bias.
% - TWO-STAGE EXPERIMENTAL STRATEGY:
%   1. Single-Factor Isolated Testing: Test each optimization module (M1, M2, M3, M4) separately against M0 Baseline to quantify individual impact.
%   2. Combinatorial Stacking & Interaction Analysis: Evaluate pairwise and multi-module combinations up to full stack (M1+M2+M3+M4) to identify Pareto frontier and systemic trade-offs (RQ4).
% - TODO: Add Factorial / Experimental Design Matrix Table here (e.g., C0: Baseline, C1: M1, C2: M2, ..., C15: M1+M2+M3+M4 full stack).

\subsection{System Implementation and Models}
\label{subsec:sys_implementation}

The experimental system is based on the enterprise-standard software stack. The system is implemented using Java versions 21 LTS with Spring Boot 4.0.0 and Spring AI version 2.0.0. The database is implemented using PostgreSQL version 16 and the pgvector extension version 0.8.0. The software stack is deployed on Google Cloud Platform, and the database is hosted on Google CloudSQL to mimic a realistic enterprise system. Northen europe datacenters are used to decrease potential delays caused by the RRT. The base language model used in the stack is Google Gemini 3.7-Flash, and in model cascading, Flash-3.5-lite and 3.1-pro are also used. In embeddings, gemini-embedding-2 with a dimension output of 768 is used to balance between efficiency and storage needs. In the semantic re-ranking, GCP's semantic-ranker-default-004 model is used 

% DRAFT NOTES - LOCKED SYSTEM IMPLEMENTATION DECISIONS:
%
% TODO: Add Visual System Architecture Diagram (figure illustrating Spring AI orchestrator, pgvector vector store, GCP Vertex Ranking API, GCP Context Cache, and 3 Gemini model tiers).
%
% 1. SOFTWARE STACK & VERSION BASELINE:
%    - Java Baseline: Java 21 LTS.
%    - Framework: Spring Boot 4.0 + Spring AI 2.0.0 (spring-ai-starter-model-google-genai).
%    - Database Engine: PostgreSQL 16 + pgvector extension (v0.8.x).
%    - Instrumentation: Spring Boot Actuator with Micrometer export for microsecond TTFT, end-to-end latency, and token accounting (RQ1 & RQ2).
%
% 2. CLOUD INFRASTRUCTURE & API GATEWAY (3 GEMINI MODEL TIERS):
%    - Infrastructure: GCP Vertex AI enterprise gateway (aiplatform.googleapis.com in us-central1 via IAM/ADC), with Google AI Studio API key fallback (generativelanguage.googleapis.com) for local testing.
%    - Generative Model Tiers (Module M4 Model Cascading):
%      1. gemini-3.5-flash-lite : Ultra-fast / low-cost tier for query routing & low-complexity single-hop tasks.
%      2. gemini-3.6-flash      : Primary workhorse for medium-complexity generation, coding, and fast context synthesis.
%      3. gemini-3.1-pro        : Flagship reasoning model for complex multi-hop synthesis (3-hop / 4-hop queries) and fallback.
%
% 3. NATIVE EMBEDDINGS ARCHITECTURE:
%    - Primary Model: GCP Native gemini-embedding-002 (2,048 token window).
%    - Dimensionality Optimization: Matryoshka Representation Learning (MRL) configured to output 768 dimensions (via output_dimensionality=768) to minimize storage overhead and vector search latency without accuracy loss.
%
% 4. TWO-STAGE RE-RANKING ARCHITECTURE:
%    - Production Standard: GCP Vertex AI Ranking API (Semantic Ranker: semantic-ranker-default-004 & semantic-ranker-fast-004).
%    - Local Development Option: Optional local fallback using TEI-hosted open cross-encoders (e.g. deberta-v3-large / jina-reranker-v2) for offline testing without cloud credentials.
%    - Workflow: Stateless cross-encoder joint self-attention re-sorting top-200 candidate passages retrieved by first-stage vector search.
%
% 5. DATABASE DEPLOYMENT & LOCAL-TO-CLOUD PORTABILITY:
%    - Local Development: PostgreSQL 16 + pgvector running in Docker container.
%    - GCP Cloud Production: Cloud SQL for PostgreSQL 16 with pgvector extension.
%    - Portability & Migration: 100% feature and SQL operator parity (<=> Cosine distance, HNSW indexing). Data migration via standard pg_dump -Fc locally and pg_restore to Cloud SQL.

\subsection{Optimization Pipeline}
\label{subsec:opt_pipeline}

There are dozens of potential ways to implement each of the optimization methods. For example, the number of tokens used in chunking could differ significantly between the methods and research papers on the methods. In addition to that, underlying model capabilities and sizes, not to mention different parameter combinations, increase the potential research space of the methods. In this subsection, the chosen methodology is presented in Table~\ref{tab:optimization_modules}. The chosen parameters are based balancing between on the occurrences in the articles in Section~\ref {sec:literature-review} and the ease of implementation and testing.

    \begin{table}[htbp]
    \centering
    \small
    \caption{Architectural specifications and technical implementation of the evaluated RAG optimization modules.}
    \label{tab:optimization_modules}
    \begin{tabularx}{\textwidth}{l p{3.2cm} X p{3.6cm}}
    \toprule
    \textbf{} & \textbf{Optimization Strategy} & \textbf{Technical Implementation and Configuration} & \textbf{Key
  References} \\
    \midrule
    \textbf{$M_0$} & Naive RAG (Baseline)
      & Fixed 512-token chunks (50-token overlap), standard dense bi-encoder retrieval ($k=5$), no re-ranking or
  caching.
      & \citep{wang2024searching, smigielski2026chunking, lewis2020rag, karpukhin2020dense} \\
    \addlinespace
    \textbf{$M_1$} & Small-to-Big Retrieval
      & Hierarchical windowing: 128-token child chunks indexed for vector retrieval, expanding to 1024-token parent
  contexts via Spring AI metadata mapping.
      & \citep{demourajunior2026adaptive, li2026bridgerag, chen2026logicrag} \\
    \addlinespace
    \textbf{$M_2$} & Two-Stage Semantic Re-ranking
      & Production-grade GCP Vertex AI Ranking API cross-encoder scoring to re-rank candidate passages and filter
  irrelevant context.
      & \citep{moreira2024enhancing, cheng2026enhancing, yu2024rankrag, wang2025infogain} \\
    \addlinespace
    \textbf{$M_3$} & Dual-Tier Prompt Caching
      & GCP Gemini Context Caching for static system instructions/contexts ($>10\text{k}/32\text{k}$ tokens) combined
  with dynamic KV prefix reuse.
      & \citep{gim2024promptcache, jin2025ragcache, lu2025turborag, yao2025cacheblend, wang2026fusionrag} \\
    \addlinespace
    \textbf{$M_4$} & Hybrid Model Cascading
      & Native Spring AI ONNX intent classifier routing low-complexity/2-hop queries to Gemini Flash and 3/4-hop
  queries to Gemini Pro, with confidence fallback.
      & \citep{chen2023frugalgpt, ong2025routellm, zhang2025ragrouter, jeong2024adaptiverag, guo2026deepsieve,
  park2025stoprag} \\
    \bottomrule
    \end{tabularx}
    \end{table}
% DRAFT NOTES:
% - Baseline (M0): Naive flat RAG (fixed 512-token chunks with 50-token overlap, bi-encoder top-5, no caching) \cite{wang2024searching, smigielski2026chunking, lewis2020rag, karpukhin2020dense}.
% - Module 1 (M1): Small-to-Big Retrieval (128 token child chunks expanding to 1024 parent windows via native Java metadata mapping in Spring AI) \cite{demourajunior2026adaptive, li2026bridgerag, chen2026logicrag}.
% - Module 2 (M2): Two-Stage Semantic Re-ranking (GCP Vertex AI Ranking API in production) \cite{moreira2024enhancing, cheng2026enhancing, yu2024rankrag, wang2025infogain}.
% - Module 3 (M3): Prompt Caching (Dual-tier: GCP Gemini Context Caching for static system context >10k/32k tokens + dynamic KV prefix reuse) \cite{gim2024promptcache, jin2025ragcache, lu2025turborag, yao2025cacheblend, wang2026fusionrag}.
% - Module 4 (M4): Hybrid Model Routing & Fallback Cascading (Native Spring AI ONNX/intent classifier routing 2-hop to Gemini Flash & 3/4-hop to Gemini Pro, with confidence fallback) \cite{chen2023frugalgpt, ong2025routellm, zhang2025ragrouter, jeong2024adaptiverag, guo2026deepsieve, park2025stoprag}.
% - Extension: Thinking Effort Routing / Reasoning Budget Control (e.g. arXiv:2603.07915, OpenReview:W34UyCRQel).

\subsection{Benchmark Datasets}
\label{subsec:bench_datasets}

% DRAFT NOTES:
% - Main Academic Benchmark: MuSiQue dataset (Trivedi et al., 2022) - sampled 100 queries (35% 2-hop, 65% 3-hop/4-hop).
% - Why MuSiQue: Eliminates single-hop reasoning shortcuts present in HotpotQA; 20 paragraphs/query (~2.5k-4k tokens) ideal for testing M1 (Small-to-Big), M2 (Re-ranking), and M3 (Prompt Caching).
% - Hop-count difficulty tiers (2-hop vs 3/4-hop) natively replace RouterBench for M4 (Model Cascading: Gemini Flash vs Pro routing).
% - Secondary Domain Corpus: Enterprise IT Customer Support Knowledge Base (500+ KB chunks) + 200 synthetic QA pairs (40% single-hop, 40% multi-hop, 20% out-of-domain bypass).


\subsection{Evaluation Setup and Metrics}
\label{subsec:eval_metrics}

% DRAFT NOTES - LOCKED EVALUATION & MEASUREMENT DECISIONS:
%
% 1. EFFICIENCY & LATENCY INSTRUMENTATION (RQ1):
%    - Metrics: Time to First Token (TTFT in ms), End-to-End Latency (ms), and micro-phase breakdown (Vector Retrieval, Vertex Re-ranking, Prompt Caching lookup, LLM TTFT, Stream Completion).
%    - Measurement Protocol: Spring Boot Actuator + Micrometer Timers. 5 iterations per query configuration (1 warm-up run discarded to eliminate JVM JIT / TLS cold-start bias; 4 recorded runs).
%    - Reporting: Report both Median (P50) and P95 tail latency for enterprise SLA analysis.
%
% 2. EXACT OPERATIONAL COST ACCOUNTING (RQ2):
%    - Execution Cost Formula: Cost ($) = (N_uncached * P_prompt) + (N_cached * P_cached) + (N_completion * P_completion) + (N_rerank_calls * P_rerank).
%    - Pricing Baseline: GCP Vertex AI official pricing tiers for gemini-3.5-flash-lite, gemini-3.6-flash, gemini-3.1-pro, Vertex Ranking API (semantic-ranker), and GCP Context Caching (75% discount on cached input tokens). Excludes fixed database storage.
%    - TODO: Write out formal LaTeX equation environment for Cost ($) during prose writing.
%
% 3. QUALITY METRICS & LLM-AS-A-JUDGE SETUP (RQ3):
%    - Framework: Offline Python RAGAS evaluation suite (Es et al., 2024).
%    - Quality Triad: Context Precision (retrieval signal-to-noise ratio), Context Recall (ground-truth coverage), Faithfulness (hallucination check against retrieved context), and Answer Relevance (semantic alignment with user query).
%    - Primary Judge Configuration: gemini-3.1-pro at T=0.0 as primary evaluation judge.
%    - THREATS TO VALIDITY & JUDGE FLEXIBILITY NOTE: Address potential LLM self-preference bias (vendor affinity). Explicitly state that judge model can be cross-validated or substituted with alternative judge LLMs (e.g. GPT-4o, Claude 3.5 Sonnet, or open-weights evaluation judges) for robustness verification.
%
% 4. TRADE-OFF EQUILIBRIUM RATIOS (RQ4):
%    - Trade-off Indicators: Practical Efficiency Ratios quantifying optimization ROI:
%      a) Quality per Dollar (QPD = RAGAS Composite Score / Total Request Cost in $)
%      b) Quality per Second (QPS = RAGAS Composite Score / End-to-End Latency in seconds)
%    - Literature Sources for Math Equations: RAGAS (Es et al., 2024; Wang et al., 2024), FrugalGPT cost/routing (Chen et al., 2023), FlashRAG benchmarking (Jin et al., 2025), FusionRAG prompt caching (Wang et al., 2026).
%    - TODO: Formulate formal LaTeX equation environments for QPD and QPS during prose writing.